problem:  
Let \((\mathbb{H}^n,\Delta,g_0)\) denote the standard Heisenberg group, with horizontal distribution \(\Delta\), canonical sub-Riemannian metric \(g_0\), and Haar measure \(\mu\).  
For a smooth \(\varphi:\mathbb{H}^n\to(0,\infty)\) that is bounded above and below by positive constants, define the conformal horizontal deformation  
\[
\langle X,Y\rangle_\varphi = \varphi^2(p)\,\langle X,Y\rangle_{g_0},\quad X,Y\in \Delta_p,
\]
and let \(d_\varphi\) denote the induced Carnot–Carathéodory distance. Consider the metric-measure structure \((\mathbb{H}^n,d_\varphi,\mu)\).  

**Question:**  
For which smooth and bounded \(\varphi\) does the deformed structure \((\mathbb{H}^n,d_\varphi,\mu)\) continue to satisfy the metric-measure curvature-dimension condition \(\mathrm{MCP}(0,N)\) with finite \(N\)?  
Can the conformal horizontal deformation decrease the effective synthetic dimension \(N\) or produce strictly stronger volume-contraction behavior compared to the standard Heisenberg structure?  
In other words, characterize all such conformal horizontal deformations that preserve or improve the synthetic nonnegative-curvature behavior encoded by \(\mathrm{MCP}(0,N)\).

Why is it a "good" problem:  
This problem avoids the known impossibility of positive curvature on the noncompact Heisenberg group while delving into the subtle rigidity of its *nonnegative* synthetic curvature behavior. It directly tests whether bounded conformal deformations can influence the optimal synthetic dimension exponent \(N\) in the measure contraction property—a deep quantitative invariant of geometric measure contraction.  

The question integrates sub-Riemannian geometry, optimal transport, and analysis on metric-measure spaces, connecting conformal geometry with synthetic curvature theory. It is simple to state yet potentially groundbreaking: resolving it would clarify whether MCP rigidity is purely structural (tied to the nilpotent group law) or flexible under analytic perturbations. This bridges conformal sub-Riemannian analysis, curvature-dimension inequalities, and geometric measure theory, exemplifying the clarity, depth, and cross-field reach characteristic of a “good” research problem.